\documentclass[../../main.tex]{subfiles}
\begin{document}
\section{Riemann 流形上的积分}\label{section:ch16-integration-on-riemannian-manifolds}

在本节中，我们把积分和 Stokes 定理的理论专门应用于 Riemann 流形。

设 $(M,g)$ 是一个定向 Riemann 流形。由\cref{proposition:ch15-volume-form}，$M$ 上存在一个唯一的体积形式 $\mathrm{d}V_g$，它在每个保定向标准正交基上取值为 $1$。如果 $M$ 不带边，这个体积形式使我们能够积分 $M$ 上任意具有紧支集的连续函数 $f$：定义
\begin{align*}
\int_M f\,\mathrm{d}V_g
\end{align*}
为 $n$-形式 $f\,\mathrm{d}V_g$ 在 $M$ 上的积分。

\begin{proposition}[Riemann 密度]\label{proposition:ch16-riemannian-density}
设 $(M,g)$ 是一个定向 Riemann 流形。由上式定义的积分在保定向局部等距下不变；它与 Riemann 度量 $g$ 自然地相伴。
\end{proposition}

在局部坐标 $(x^1,\dots,x^n)$ 下，若 $g_{ij}$ 是度量系数，则
\begin{align*}
\mathrm{d}V_g=\sqrt{\det(g_{ij})}\,\mathrm{d}x^1\wedge\cdots\wedge\mathrm{d}x^n.
\end{align*}

\begin{lemma}\label{lemma:ch16-divergence-coord}
设 $X$ 是定向 Riemann 流形 $M$ 上的光滑向量场。在局部坐标下，散度 $\operatorname{div}X$ 由下式给出：
\begin{align*}
\operatorname{div}X=\frac{1}{\sqrt{\det g}}\frac{\partial}{\partial x^i}\left(\sqrt{\det g}\,X^i\right).
\end{align*}
\end{lemma}

\begin{theorem}[散度定理]\label{theorem:ch16-divergence-theorem}
设 $(M,g)$ 是一个定向 $n$ 维 Riemann 带边流形，$X$ 是 $M$ 上一个具有紧支集的光滑向量场。则
\begin{align}
\int_M(\operatorname{div}X)\,\mathrm{d}V_g=\int_{\partial M}\langle X,N\rangle\,\mathrm{d}V_{\widetilde g},
\end{align}
其中 $N$ 是 $\partial M$ 上指向外的单位法向量场，$\widetilde g$ 是 $\partial M$ 上的诱导度量。
\end{theorem}

\begin{proof}
令 $\omega=X\lrcorner\,\mathrm{d}V_g$。由 Cartan 魔公式，$\mathrm{d}\omega=(\operatorname{div}X)\,\mathrm{d}V_g$。在边界上，$\omega$ 的取值等于 $\langle X,N\rangle\,\mathrm{d}V_{\widetilde g}$。对 $\omega$ 应用 Stokes 定理即得。
\end{proof}

\begin{figure}[ht]
\centering
\includegraphics[width=0.85\textwidth]{fig-16-14.png}
\caption{同伦的分段光滑曲线段}\label{figure:ch16-homotopic-curve-segments}
\end{figure}

\begin{proposition}\label{proposition:ch16-surface-integral-vector-field}
设 $S\subseteq M$ 是一个定向嵌入超曲面，$N$ 是 $S$ 上确定定向的单位法向量场。对 $M$ 上任意光滑向量场 $X$，$X$ 穿过 $S$ 的通量定义为
\begin{align*}
\int_S\langle X,N\rangle\,\mathrm{d}V_{\widetilde g}.
\end{align*}
\end{proposition}

\begin{figure}[ht]
\centering
\includegraphics[width=0.6\textwidth]{fig-16-15.png}
\caption{散度的几何解释}\label{figure:ch16-divergence-geometric}
\end{figure}

\begin{theorem}[曲面上的 Stokes 公式]\label{theorem:ch16-stokes-surface-integrals}
设 $S\subseteq\mathbb{R}^3$ 是一个定向光滑带边曲面，$\partial S$ 带有诱导定向。对 $S$ 的邻域上任意光滑向量场 $X$，有
\begin{align*}
\int_S(\operatorname{curl}X)\cdot N\,\mathrm{d}S=\int_{\partial S}X\cdot\mathrm{d}\mathbf{s}.
\end{align*}
\end{theorem}

\begin{figure}[ht]
\centering
\includegraphics[width=0.6\textwidth]{fig-16-16.png}
\caption{曲面积分的设定}\label{figure:ch16-surface-integral-setup}
\end{figure}

\begin{exercise}\label{exercise:ch16-43}
设 $(M,g)$ 是一个紧定向 Riemann 流形（不带边）。证明 $\int_M\operatorname{div}X\,\mathrm{d}V_g=0$ 对所有光滑向量场 $X$ 成立。
\end{exercise}

\begin{exercise}\label{exercise:ch16-44}
设 $(M,g)$ 是一个紧定向 Riemann 带边流形。证明 Green 第一恒等式与 Green 第二恒等式。
\end{exercise}

\begin{theorem}[不可定向情形下的散度定理]\label{theorem:ch16-divergence-nonorientable}
即使 $M$ 不可定向，用 Riemann 密度定义的积分仍然有意义，散度定理在作相应改写后依然成立。
\end{theorem}
\end{document}
